\documentclass[11pt]{article}
\usepackage[margin=1in]{geometry}
\usepackage{amsmath,amssymb,amsthm}
\usepackage{booktabs}
% Latin Modern: scalable Type 1 outlines of Computer Modern. microtype's font expansion
% refuses to run against the bitmap CM fonts this machine's TeX resolves by default
% ("auto expansion is only possible with scalable fonts"), which is a fatal error rather
% than a warning. Same typeface, so nothing about the document's appearance changes.
\usepackage{lmodern}
\usepackage{microtype}
\usepackage[hidelinks]{hyperref}
\usepackage{enumitem}

\newtheorem{lemma}{Lemma}
\newtheorem{proposition}{Proposition}
\theoremstyle{remark}
\newtheorem{remark}{Remark}

\DeclareMathOperator*{\argmin}{arg\,min}
\DeclareMathOperator{\Cov}{Cov}
\DeclareMathOperator{\Var}{Var}
\newcommand{\R}{\mathbb{R}}
\newcommand{\1}{\mathbf{1}}
\newcommand{\Ind}[1]{\mathbb{I}\!\left[#1\right]}
\newcommand{\VT}{V_T}

\title{\textbf{\texttt{mlr+mom+K}}: a complete derivation\\[2pt]
\large A multinomial-logit fitness model over allele features and frequency momentum,\\
convolved with a consensus mutation kernel}
\author{antiGen}
\date{28 September 2026}

\begin{document}
\maketitle

\begin{abstract}
\noindent
This note derives, in full, the forecasting model referred to in the run logs as
\texttt{mlr+mom+K}. The model predicts the genotype composition of the next influenza
A/H3N2 season as a distribution over observed genotypes \emph{and their one-substitution
neighbours}. It has three parts, each derived separately below: a base measure obtained
by exponentially discounting past observations; an exponential tilt of that base measure
whose exponent is linear in one-hot allele indicators and a frequency-momentum column,
fitted by penalised multinomial likelihood against a nearest-neighbour assignment
target; and a convolution with a mutation kernel that proposes, for each genotype, the
substitutions toward alleles the population is moving \emph{to}. Every forecast is scored
by the exact earth-mover's distance to the realised season under a Hamming ground metric.
We state the estimator's convexity, the gradient used to fit it, the identification
structure of the design, the mass-conservation and nesting properties of the
convolution, and the selection protocol under which every hyperparameter is chosen ---
including the kernel's own timescale and temperature, which earlier versions of this note
fixed by an oracle diagnostic and which are now resolved per origin from information
available at that origin. We close with the empirical result and the caveats that
presently govern it.
\end{abstract}

\tableofcontents

\section{The forecasting problem}

\subsection{Data and notation}

The data are dated amino-acid sequences. Writing $\mathcal{D} = \{(t_i, g_i)\}_{i=1}^{N}$
for sequence $i$ observed at decimal time $t_i$ with genotype $g_i$, each genotype is a
string over the $566$-residue concatenated translation
$\texttt{SigPep} \mathbin\Vert \texttt{HA1} \mathbin\Vert \texttt{HA2}$ of the HA segment.
On the backbone panel $N = 1{,}683$ tips carry $1{,}063$ distinct genotypes.

A \emph{season} $s$ runs from 1 October of year $s-1$ to 1 April of year $s$. A forecast
made at origin season $s$ has \emph{horizon}
\begin{equation}
  T \;=\; s + \tfrac14,
\end{equation}
i.e.\ 1 April of season $s$, and predicts the composition of season $s+1$.

\subsection{Support, target, ground metric}

\paragraph{Support.} The forecast is a probability vector over the genotypes observable
at the horizon within a lookback window $L$:
\begin{equation}
  \VT \;=\; \bigl\{\, g_i \;:\; T - L < t_i \le T \,\bigr\},
  \qquad L = 5 \text{ years}.
\end{equation}
Typical sizes are $|\VT| \in [271, 442]$. Nothing dated after $T$ enters any quantity
below; this is the property that makes every component leak-free, and it is enforced by
test rather than by convention.

\paragraph{Target.} Let $G_{s+1}$ be the distinct genotypes of season $s+1$ and
$y \in \Delta^{|G_{s+1}|-1}$ their empirical frequencies.

\paragraph{Ground metric.} For $v \in \VT$ and $j \in G_{s+1}$,
\begin{equation}
  M_{vj} \;=\; d_{\mathrm H}(v, j) \;=\; \sum_{q=1}^{566} \Ind{v_q \neq j_q},
\end{equation}
the Hamming distance in amino-acid substitutions. All reported quantities are in these
units.

\subsection{The scoring functional}

A forecast $x \in \Delta^{|\VT|-1}$ is scored by the exact $1$-Wasserstein distance
(earth-mover's distance) to the realised season:
\begin{equation}
  W(x, y) \;=\; \min_{\Gamma \in \R_{\ge0}^{|\VT| \times |G_{s+1}|}}
  \; \sum_{v,j} \Gamma_{vj} M_{vj}
  \quad\text{s.t.}\quad \Gamma \1 = x,\;\; \Gamma^\top \1 = y .
  \label{eq:emd}
\end{equation}
This is a linear program, solved exactly (network simplex, \texttt{float64}). In the run
artifacts the value $W(x,y)$ for a family \texttt{fam} is recorded under the key
\texttt{emd\_fam}; the prefix names the \emph{metric}, not the model.

Two reference levels accompany every panel:
\begin{align}
  \text{persistence} &:\quad W(x^{\mathrm{seas}}, y),
    \quad x^{\mathrm{seas}} = \text{uniform over season-$s$ sequences},\\
  \text{floor} &:\quad \min_{x \in \Delta^{|\VT|-1}} W(x,y)
    \;=\; \sum_j y_j \min_{v \in \VT} M_{vj},
  \label{eq:floor}
\end{align}
the second because the optimal reweighting sends every target's mass to its nearest
support point. The floor is what \emph{any} reweighting of $\VT$ could achieve; it is
the correct bar for every model in this note except the kernel arm, for which
Section~\ref{sec:extfloor} derives a lower one.

\section{The base measure}

\subsection{Exponentially discounted level}

For a timescale $\tau>0$ define the discounted genotype shares
\begin{equation}
  \pi^{(\tau)}_v \;=\;
  \frac{\displaystyle\sum_{i \,:\, g_i = v,\; T-L < t_i \le T} e^{-(T-t_i)/\tau}}
       {\displaystyle\sum_{u \in \VT}\;\sum_{i \,:\, g_i = u,\; T-L < t_i \le T} e^{-(T-t_i)/\tau}},
  \qquad v \in \VT,
  \label{eq:level}
\end{equation}
with the convention $\pi^{(\infty)}_v \propto \#\{i : g_i = v\}$ over the window. A decay
rather than a hard window is used so that every genotype in $\VT$ retains strictly
positive weight at every $\tau$: recency is expressed without discarding support, which
matters because the support also determines the floor \eqref{eq:floor}.

\subsection{Uniform smoothing}

The tilt in Section~\ref{sec:tilt} acts multiplicatively on the base measure, so a
genotype seen once --- whose share is estimated from a single sequence --- would have a
badly estimated number exponentiated. The base measure is therefore mixed with the
uniform distribution:
\begin{equation}
  b_v \;=\; (1-\alpha)\,\pi^{(\tau_f)}_v \;+\; \frac{\alpha}{|\VT|},
  \qquad \alpha \in [0,1).
  \label{eq:base}
\end{equation}
This bounds below how small, and hence how poorly estimated, any base probability can be.
$\alpha$ is selected per origin (Section~\ref{sec:selection}); on the backbone panel the
selector chose $\alpha = 0$ at every origin.

\section{The design}

\subsection{Allele features}

At each of the $566$ sites, let the \emph{varying} sites be those carrying more than one
allele at panel-wide frequency $\ge \phi$ ($\phi = 0.005$). Enumerating the surviving
(site, allele) pairs gives the feature index set $\mathcal{F}$, $|\mathcal{F}| = d$, over
$103$ varying sites with $d = 239$ on the backbone panel. The one-hot block is
\begin{equation}
  S \in \{0,1\}^{|\VT| \times d},
  \qquad S_{v,(q,a)} \;=\; \Ind{v_q = a}.
\end{equation}

\begin{remark}[No intercept]
Each row of $S$ sums to the number of varying sites at which $v$ carries an in-grid
allele --- between $100$ and $103$, hence nearly but not exactly constant. An explicit
intercept would be unidentified, since adding a constant to the exponent cancels in
\eqref{eq:tilt}; see Lemma~\ref{lem:shift}. None is included.
\end{remark}

\subsection{The momentum column}

Let $\tau_f < \tau_s$ be a fast and a slow timescale. The raw log-ratio of the two
discounted levels,
\begin{equation}
  \tilde m_v \;=\; \log \pi^{(\tau_f)}_v \;-\; \log \pi^{(\tau_s)}_v,
\end{equation}
is positive for genotypes whose share is rising and negative for those whose share is
falling: it is a discrete-time estimate of $\tfrac{d}{dt}\log \pi_v$ smoothed over two
horizons. It is then centred under the base measure,
\begin{equation}
  m_v \;=\; \tilde m_v \;-\; \sum_{u \in \VT} \pi^{(\tau_f)}_u \, \tilde m_u .
  \label{eq:mom}
\end{equation}
The design is the concatenation
\begin{equation}
  Z \;=\; \bigl[\, S \;\big\vert\; m \,\bigr] \in \R^{|\VT| \times (d+1)},
  \qquad \theta = (\beta, \gamma) \in \R^{d} \times \R .
\end{equation}

\begin{lemma}[Shift invariance; why centring is free]\label{lem:shift}
Let $x(\theta)$ be as in \eqref{eq:tilt}. For any $c \in \R$, replacing the column $m$ by
$m - c\1$ leaves $x(\theta)$ unchanged for every $\theta$.
\end{lemma}
\begin{proof}
Replacing $m$ by $m - c\1$ replaces the exponent $z_v^\top\theta$ by
$z_v^\top\theta - c\gamma$, a constant in $v$. Then
$b_v e^{z_v^\top\theta - c\gamma} = e^{-c\gamma} b_v e^{z_v^\top\theta}$, and the common
factor $e^{-c\gamma}$ cancels between numerator and denominator of \eqref{eq:tilt}.
\end{proof}

Centring therefore changes no prediction. It is done so that the penalty on $\gamma$
(Section~\ref{sec:penalty}) means ``how much tilt'' rather than ``how much tilt plus an
offset'', which would otherwise make the regularisation path depend on an arbitrary
additive constant.

\section{The forecast}\label{sec:tilt}

The model is an exponential tilt --- equivalently a conditional logit with an offset ---
of the base measure:
\begin{equation}
  \boxed{\;
  x_v(\theta) \;=\;
  \frac{b_v \, \exp\!\bigl(\beta^\top S_v + \gamma\, m_v\bigr)}
       {\displaystyle\sum_{u \in \VT} b_u \, \exp\!\bigl(\beta^\top S_u + \gamma\, m_u\bigr)},
  \qquad v \in \VT .
  \;}
  \label{eq:tilt}
\end{equation}
Writing $f = Z\theta$, the implementation evaluates
$x \propto b \odot \exp(f - \max_v f_v)$, which is the standard shift for numerical
stability and is exact by Lemma~\ref{lem:shift}.

Interpretation: $\exp(\beta^\top S_v + \gamma m_v)$ is a relative fitness. A genotype's
predicted share is its current (discounted) share multiplied by a fitness that depends on
which alleles it carries and on whether its own frequency trajectory is rising. This is
the Malthusian form $x_v \propto b_v e^{f_v}$ used throughout the field; what is specific
here is the feature set and the estimator.

\section{The estimator}

\subsection{Why not minimise the score directly}

The natural estimator would minimise \eqref{eq:emd} in $\theta$. We do not, for two
reasons. First, $\theta \mapsto W(x(\theta), y)$ is not convex: $W$ is convex in $x$, but
$x(\theta)$ is not an affine function of $\theta$, and the composition need not be convex.
Second, each objective evaluation costs a transport solve. The estimator below is convex,
has a closed-form gradient, and costs no transport solve at all; the transport LP is paid
once per scored forecast, at evaluation time. \emph{Everything reported is still scored by
\eqref{eq:emd}} --- the surrogate is used for fitting only.

\subsection{The assignment target}

Exact genotypes rarely recur: on this panel only about $6\%$ of a season's sequences carry
a genotype already seen before 1 April. A multinomial likelihood over exact genotype
identity would therefore discard $94\%$ of the data. Instead, each next-season sequence is
assigned to its nearest already-observed genotype, ties split equally. With
$\mu_j = \min_{u \in \VT} M_{uj}$,
\begin{equation}
  q_v \;=\; \sum_{j \in G_{s+1}} y_j \cdot
  \frac{\Ind{M_{vj} = \mu_j}}{\bigl|\{u \in \VT : M_{uj} = \mu_j\}\bigr|},
  \qquad \sum_{v \in \VT} q_v = 1 .
  \label{eq:assign}
\end{equation}
This is precisely the transport plan attaining the floor \eqref{eq:floor}. It is a
\emph{surrogate target for estimation}; it is never a scoring target.

\subsection{Penalised multinomial cross-entropy}

Let $\mathcal{T}_s$ be the training transitions available at origin $s$
(Section~\ref{sec:selection}). For $r \in \mathcal{T}_s$ write $Z^{(r)}, b^{(r)},
q^{(r)}$ for that transition's design, base measure and assignment target, and
$x^{(r)}(\theta)$ for \eqref{eq:tilt} evaluated on it. The estimator is
\begin{equation}
  \boxed{\;
  \hat\theta \;=\; \argmin_{\theta \in \R^{d+1}}\;
  \underbrace{\frac{1}{|\mathcal{T}_s|}\sum_{r \in \mathcal{T}_s}
  \Bigl(-\sum_{v} q^{(r)}_v \log x^{(r)}_v(\theta)\Bigr)}_{\textstyle \ell(\theta)}
  \;+\; \theta^\top P \theta ,
  \;}
  \label{eq:est}
\end{equation}
with $P$ the diagonal penalty of Section~\ref{sec:penalty}.

\begin{proposition}[Gradient]\label{prop:grad}
For a single transition with design $Z$, base $b$, target $q$ and total target mass
$q_{\mathrm{tot}} = \sum_v q_v$,
\begin{equation}
  \nabla_\theta \Bigl(-\sum_v q_v \log x_v(\theta)\Bigr)
  \;=\; Z^\top\bigl(q_{\mathrm{tot}}\, x(\theta) - q\bigr).
  \label{eq:grad}
\end{equation}
\end{proposition}
\begin{proof}
Write $f = Z\theta$, $w_v = b_v e^{f_v}$, $A = \sum_u w_u$, $x = w/A$. Then
$\log x_v = \log b_v + z_v^\top\theta - \log A$ and
\[
  \nabla_\theta \log A
  = \frac{1}{A}\sum_u w_u z_u = \sum_u x_u z_u = Z^\top x .
\]
Hence $\nabla_\theta \log x_v = z_v - Z^\top x$, and
\[
  \nabla_\theta\Bigl(-\sum_v q_v \log x_v\Bigr)
  = -\sum_v q_v\bigl(z_v - Z^\top x\bigr)
  = -Z^\top q + q_{\mathrm{tot}} Z^\top x . \qedhere
\]
\end{proof}

This is the textbook multinomial-logit gradient, ``observed minus expected''. The factor
$q_{\mathrm{tot}}$ is inert here, since \eqref{eq:assign} normalises to one; it is
retained because the same routine serves a within-season fit whose bundles carry unnormalised
shares, and omitting it there would make the loss and its gradient different functions.

\begin{proposition}[Convexity]\label{prop:convex}
$\ell$ is convex in $\theta$, with Hessian for a single transition
\begin{equation}
  \nabla^2_\theta \Bigl(-\sum_v q_v \log x_v\Bigr)
  \;=\; q_{\mathrm{tot}}\; Z^\top\bigl(\operatorname{diag}(x) - x x^\top\bigr) Z \;\succeq\; 0 .
\end{equation}
\end{proposition}
\begin{proof}
Differentiating \eqref{eq:grad}, $\nabla_\theta x = (\operatorname{diag}(x) - xx^\top)Z$.
For any $u$, $u^\top Z^\top(\operatorname{diag}(x) - xx^\top)Zu = \Var_{v\sim x}(z_v^\top u) \ge 0$,
the variance of a real random variable under the probability vector $x$. A nonnegative
combination over $r \in \mathcal{T}_s$ preserves positive semidefiniteness, and
$\theta^\top P\theta$ is convex, so \eqref{eq:est} is a convex program.
\end{proof}

Consequently the fit has a unique optimal value, and the solution does not depend on
initialisation. Optimisation is by L-BFGS-B with analytic gradient \eqref{eq:grad}.

\subsection{Anisotropic penalty}\label{sec:penalty}

The penalty is diagonal but not isotropic:
\begin{equation}
  P \;=\; \operatorname{diag}\bigl(\underbrace{\lambda,\dots,\lambda}_{d
  \text{ allele features}},\; \gamma_{\mathrm{pen}}\bigr),
  \qquad
  \theta^\top P \theta \;=\; \lambda\|\beta\|_2^2 + \gamma_{\mathrm{pen}}\gamma^2 .
\end{equation}
The rationale is that the two blocks are doing different jobs. The $d = 239$ allele
coefficients are many relative to the two or three training transitions available and
must be shrunk. The single momentum coefficient \emph{is} the term the family is named
for; penalising it at $\lambda$ would switch it off by default. It nevertheless carries a
small penalty of its own, so that a single training transition cannot drive it away.

\begin{remark}[Preconditioning]
The implementation substitutes $\theta = P^{-1/2} u$ and minimises
$\ell(P^{-1/2}u) + \|u\|_2^2$. This is an exact reparameterisation that whitens the
quadratic term; it changes the iterates of L-BFGS, not the optimum.
\end{remark}

\section{The mutation kernel}\label{sec:kernel}

Everything so far redistributes mass among \emph{observed} genotypes. The kernel extends
the forecast to genotypes one substitution away from an observed one.

\subsection{The candidate set}

Let $Q$ be the varying sites and, for each $q \in Q$, let $\mathcal{A}_q$ be the residues
\emph{observed in the support} at site $q$. The candidate columns are
\begin{equation}
  \mathcal{C} \;=\; \bigl\{(q,b) \;:\; q \in Q,\; b \in \mathcal{A}_q \bigr\},
\end{equation}
and for a parent $v \in \VT$ the candidate $v^{q\to b}$ is $v$ with site $q$ set to $b$,
defined for $b \neq v_q$. Sites carrying a single residue contribute nothing, since
substituting to the residue already present is a no-op. Restricting to observed residues
is what keeps the arm leak-free; it is also why the extended floor below is $2.780$ rather
than the $2.692$ of the full Hamming ball.

\subsection{Site--allele frequencies and their extrapolation}

Define, for timescale $\tau$ and site $q$,
\begin{equation}
  f^{(\tau)}_q(b) \;=\;
  \frac{\sum_{i\,:\, t_i \le T,\; (g_i)_q = b} e^{-(T-t_i)/\tau}}
       {\sum_{i\,:\, t_i \le T} e^{-(T-t_i)/\tau}},
  \qquad \sum_{b} f^{(\tau)}_q(b) = 1 .
  \label{eq:saf}
\end{equation}
Only sequences dated at or before $T$ contribute. With a fast timescale $\tau_e$ and a
slow one $3\tau_e$, the kernel proposes toward
\begin{equation}
  \log \hat f_q(b) \;\;\propto\;\;
  \log f^{(\tau_e)}_q(b) \;+\; \delta\Bigl(
  \log \max\bigl(f^{(\tau_e)}_q(b), \phi\bigr) -
  \log \max\bigl(f^{(3\tau_e)}_q(b), \phi\bigr)\Bigr),
  \label{eq:fhat}
\end{equation}
renormalised per site. The fast timescale $\tau_e$ is \emph{selected per origin}
(\S\ref{sec:selection}); the ratio $3$ between the slow and fast windows is declared, not
selected, and is borrowed from the panel's own $\tau_s$ grid. Here $\delta = 1$ for
\texttt{kcmom} and $\delta = 0$ for \texttt{kcons}, which is the matched control: \texttt{kcons} proposes toward where the
population \emph{is}, \texttt{kcmom} toward where it is \emph{heading}. Equation
\eqref{eq:fhat} is the site--allele reading of the genotype-level momentum column
\eqref{eq:mom}.

The floor $\phi$ (the panel's own minimum allele frequency, $0.005$) is applied
\emph{inside the ratio only}; the level term keeps its exact share. Without it, an allele
absent from the slow window but present in the fast one yields an infinite ratio and the
arm would place a site's entire proposal mass on whichever residue happened to appear most
recently.

\subsection{The proposal}

The proposal factorises into a site propensity and a residue propensity,
$\pi(q,b \mid v) \propto s_q(v)\,\psi_q(v_q \to b)$, with
\begin{equation}
  \boxed{\;
  \pi(q,b \mid v) \;\propto\;
  \underbrace{\hat f_q(b)}_{\psi}\;\cdot\;
  \underbrace{\bigl(1 - \hat f_q(v_q)\bigr)}_{s},
  \qquad b \neq v_q .
  \;}
  \label{eq:prop}
\end{equation}
The first factor proposes the allele the population is moving toward. The second is the
entire conditioning on $v$: if $v$ already carries the rising allele at site $q$ then
$1 - \hat f_q(v_q)$ is small and little proposal mass is spent there. The kernel pushes
$v$ only at the sites where $v$'s own residue is the one being left behind.

Writing $L_v(q,b) = \log \hat f_q(b) + \log\bigl(1 - \hat f_q(v_q)\bigr)$ and
$L_v(q,b) = -\infty$ when $b = v_q$, the realised proposal is a tempered, optionally
truncated softmax,
\begin{equation}
  \pi(\cdot \mid v) \;=\; \operatorname{softmax}\bigl(L_v / \tau_k\bigr)
  \quad\text{restricted to its top } m \text{ columns},
  \label{eq:softmax}
\end{equation}
with $m = 0$ meaning no truncation. As $\tau_k \to 0$ the proposal approaches a point mass
on the single most favourable substitution per parent. Both $\tau_k$ and $m$ are selected
per origin (\S\ref{sec:selection}).

It is worth recording what $\tau_k$ does and does not control, because the selected values
are small enough to look alarming. Between $\tau_k = 0.02$ and $\tau_k = 0.001$ the number
of materialised candidates falls by a factor of roughly fourteen, while the mass-weighted
descent statistics move in the third decimal ($p_{\mathrm{down}}$ $0.7425 \to 0.7415$,
$p_{\mathrm{flat}}$ $0.0684 \to 0.0675$ at $\tau_e = 0.25$). The discarded candidates carry
no mass: the small-$\tau_k$ proposal is the same distribution with a numerically negligible
tail pruned, not a qualitatively sharper kernel.

\subsection{Convolution}

Let $K$ be the operator sending mass at $v$ to $\{v^{q\to b}\}$ in proportions
\eqref{eq:softmax}, and let $p_v \in \{0,1\}$ indicate whether $v$ has at least one
candidate. The forecast becomes
\begin{equation}
  \boxed{\; x' \;=\; (1-\mu)\,x \;+\; \mu\, K x , \qquad \mu \in [0,1]. \;}
  \label{eq:conv}
\end{equation}
Operationally the implementation keeps $x_v(1 - \mu p_v)$ on the original support and
emits $\mu \sum_v x_v p_v$ onto the candidate block, which is \eqref{eq:conv} exactly.

\begin{lemma}[Mass conservation]
$\sum x' = \sum x$ for every $\mu \in [0,1]$, including for parents with no candidate.
\end{lemma}
\begin{proof}
The retained mass is $\sum_v x_v(1 - \mu p_v)$ and the emitted mass is
$\mu\sum_v x_v p_v$, since $K$ is row-stochastic on $\{v : p_v = 1\}$. The two sum to
$\sum_v x_v$. A parent with $p_v = 0$ has an empty row of $K$, so $\mu$ multiplies nothing
and its mass is retained in full.
\end{proof}

\begin{lemma}[Nesting]\label{lem:nest}
At $\mu = 0$, $x' = x$ bit-identically, with no candidate block materialised.
\end{lemma}

Lemma~\ref{lem:nest} is what makes the kernel family \emph{selectable}: it contains its
own parent as the $\mu = 0$ member, so the choice between \texttt{mlr+mom} and
\texttt{mlr+mom+K} is a choice of $\mu$ within one family, made by the ordinary selector
on past data alone. It is also the identity used to verify that two independent runs are
paired --- the parent columns of the kernel run must reproduce the parent run exactly, and
do, to $0$.

\subsection{Cost block and the extended floor}\label{sec:extfloor}

Scoring \eqref{eq:emd} on the extended support requires the distance from each candidate
to each target. No candidate sequence is materialised: since $v^{q\to b}$ differs from $v$
at exactly one site,
\begin{equation}
  M_{v^{q\to b},\, j} \;=\; M_{vj} \;-\; \Ind{v_q = j_q} \;+\; \Ind{b \neq j_q},
  \label{eq:costrec}
\end{equation}
so the block is obtained from the parent rows by a rank-one correction per column. This
recursion counts agreeing residues and is therefore specific to the Hamming metric; the
implementation raises rather than silently falling back under another ground metric.

Because the forecast may now place mass outside $\VT$, the bar in \eqref{eq:floor} is no
longer the right one. The relevant floor is
\begin{equation}
  \text{floor}^{\mathrm{ext}}
  \;=\; \sum_j y_j \min_{u \,\in\, \VT \cup \mathcal{N}_1(\VT)} M_{uj},
\end{equation}
where $\mathcal{N}_1$ is the one-substitution neighbourhood over observed residues. On the
panel reported in \S\ref{sec:result} $\text{floor} = 3.389$ while
$\text{floor}^{\mathrm{ext}} = 2.541$; on the earlier published panel the pair is $3.632$
and $2.780$. The
earth-mover's distances remain directly comparable across all models --- same targets,
same cost matrix --- but any statement of the form ``fraction of the achievable gain
captured'' must use each model's own floor.

\section{Selection protocol}\label{sec:selection}

\subsection{What is chosen per origin, and on what}

The hyperparameters $(\tau_f, \tau_s, \lambda, W, \alpha, \gamma_{\mathrm{pen}}, \mu)$
are re-chosen at every origin. Let $\mathcal{O} = (s_1 < s_2 < \dots)$ be the ordered
transitions and let origin $s_i$ be the one being predicted. Then:
\begin{enumerate}[noitemsep]
  \item the \emph{training set} is the $W$ most recent transitions strictly before $s_i$
        ($W = 0$ meaning all of them);
  \item the \emph{validation folds} are the $k = 3$ most recent transitions before $s_i$,
        each scored by a fit that saw only what preceded \emph{it};
  \item each grid combination is scored by mean EMD over those folds;
  \item the forecasts of the best $k' = 5$ combinations are averaged.
\end{enumerate}
Recent folds rather than all of them are used because the record is not stationary. Grid
collapse is applied per family: fields a family cannot see are not swept, so that the
argmin is never decided by ties among identical predictions.

\subsection{The kernel's own constants}\label{sec:gatesel}

The triple $(\tau_e, \tau_k, m)$ cannot ride in the grid above, because the grid is scored
by forecast EMD on the validation folds and each additional field multiplies the cost of
an already expensive search. It is chosen instead by a \emph{descent gate}. Writing
$\varphi_u = \min_j (M_{uj} - z^\star_j)$ for the $c$-transform of the optimal target
potential of the unextended transport, Danskin's theorem gives the right-derivative of the
convolved score at $\mu = 0$ as
\begin{equation}
  \frac{\partial W_1}{\partial \mu}\Bigl|_{0}
  \;=\; \sum_v x_v \sum_{(q,b)} \pi(q,b \mid v)\,
        \bigl[\varphi_{v^{q\to b}} - \varphi_v\bigr]
  \;=\; 1 - p_{\mathrm{flat}} - 2p_{\mathrm{down}},
  \label{eq:deriv}
\end{equation}
every bracket being exactly $-1$, $0$ or $+1$ because the Hamming cost is integer-valued.
This is cheap --- one transport solve per transition, no forecast scored --- but it is an
\textbf{oracle} quantity: $z^\star$ reads where the season actually went.

That is the whole difficulty, and it is why this subsection exists. Measuring
\eqref{eq:deriv} on a transition and then \emph{scoring} that transition with the
resulting constants is selection on the evaluation set. The protocol is therefore:

\begin{quote}
Gate row $j$ reads the targets of season $j+1$, which close at $(j+1)+\tfrac14$. Origin
$s_i$ forecasts from $s_i + \tfrac14$. Row $j$ is readable at origin $s_i$ exactly when
$(j+1)+\tfrac14 \le s_i + \tfrac14$, i.e.\ $j < s_i$. At each origin, $(\tau_e, \tau_k, m)$
is the argmin of \eqref{eq:deriv} averaged over the readable rows, restricted to
configurations measured on all of them.
\end{quote}

This is the same ``target has closed'' rule the training window and the validation folds
obey, stated on the gate's index rather than the forecast's. A configuration missing from
any readable transition is dropped rather than averaged over the subset it has, since a
mean over four transitions and a mean over five are not comparable. An origin with no
readable row raises rather than falling back to a declared default, which would leave part
of a panel causal and part of it not with nothing in the output saying which.

\medskip
\noindent
With that, \textbf{every quantity entering the forecast at $s_i$ --- coefficients,
hyperparameters, and the kernel's own constants --- is a function of data strictly earlier
than $s_i$.} Earlier versions of this note could not make that claim: they fixed
$(\tau_e, \tau_k, m) = (0.25, 0.02, 0)$ by pooling \eqref{eq:deriv} over transitions
$\{2017, 2018, 2020, 2021, 2022\}$, five of which are also scored. \S\ref{sec:caveats}
records what the correction cost.

The descent surface is flat where it matters, which is why it can be resolved from short
prefixes without damage: averaged over all eleven transitions, $\tau_e$ has an interior
optimum at $0.5$ ($-0.5602$, against $-0.5504$ at $0.25$ and $-0.5416$ at $1$) and
$\tau_k$ is a plateau from $0.001$ to $0.02$ with a spread of $0.005$. The per-origin
choices wander inside that plateau rather than disagreeing with each other.

\subsection{Averaging}

For predictions $x^{(1)},\dots,x^{(k')}$ possibly living on different extended supports,
the average is taken in the extended space: the support parts average pointwise and every
candidate block is carried along with its mass divided by $k'$. Blocks from different
combinations are never merged, since their candidate indexings are not comparable.

\section{Result}\label{sec:result}

On the backbone panel with the evaluation masked on season 2021 (eight scored
transitions; nothing held out of any fit), with $\delta = 1$ and $(\tau_e, \tau_k, m)$
resolved per origin by \S\ref{sec:gatesel}. The comparison models are given the same
per-origin selector, on the same folds with the same averaging, rather than being held at
the constants their papers ship; \S\ref{sec:caveats} says what that is worth and in which
direction.

\begin{center}
\begin{tabular}{lrrr}
\toprule
model & mean EMD (AA) & vs.\ persistence & transitions won \\
\midrule
\texttt{mlr+mom+K} & \textbf{7.488} & $+6.11\%$ & $8/8$ \\
\texttt{distance\_from\_consensus} & 7.566 & $+5.21\%$ & $6/8$ \\
\texttt{ne\_star-lbi} & 7.624 & $+4.80\%$ & $6/8$ \\
\texttt{mlr+mom} & 7.661 & $+3.86\%$ & $8/8$ \\
\texttt{lbi} & 7.677 & $+4.20\%$ & $5/8$ \\
\texttt{ne\_star} (Huddleston 2020) & 7.772 & $+2.35\%$ & $6/8$ \\
\midrule
persistence & 7.970 & --- & --- \\
floor (reweighting) & 3.389 & --- & --- \\
floor (extended support) & 2.541 & --- & --- \\
\bottomrule
\end{tabular}
\end{center}

\medskip
\noindent
``Transitions won'' counts wins against persistence, which is the column all rows share.
The paired gain of the kernel over its own parent is $-0.173$ AA with a $95\%$ bootstrap
interval of $[-0.272, -0.084]$, improving $7$ of $8$ transitions --- the exception being
$2019 \to 2020$, at $+0.014$ AA. So \texttt{mlr+mom+K} beats persistence everywhere and
beats its own parent almost everywhere; the two counts differ and both are reported.

The selector chose $\mu \in \{0.05, 0.1, 0.2, 0.4\}$ across the scored origins, and
$\mu = 0$ at the masked origin $2020$, where by Lemma~\ref{lem:nest} it returns the parent
bit-identically. Nothing above is separable from the rows beneath it: every paired interval
between the top five models straddles zero. What is not shared is the win count --- $8/8$
belongs to \texttt{mlr+mom+K} and \texttt{mlr+mom} alone.

\medskip
\noindent
\textbf{Out of panel.} On the 23 validation origins of the Huddleston 2020 panel
(1996--2014 sequence data, a different corpus and a six-month pivot grid), the same kernel
takes \texttt{mlr+mom} from $5.896$ to $5.750$ against a naive baseline of $6.370$: a
paired gain of $-0.1455$ AA, $95\%$ interval $[-0.206, -0.079]$, improving $21$ of $23$
origins. The four families shared with the reference run on that panel reproduce it to
$0$, so the pairing is exact. The eight test origins of that panel are not yet scored.

\section{Caveats}\label{sec:caveats}

\begin{enumerate}
\item \textbf{What the causal correction cost, and what it did not settle.} The previous
version of this note fixed $(\tau_e, \tau_k, m)$ by pooling the oracle diagnostic
\eqref{eq:deriv} over five of the eight scored transitions, and reported $7.476$ /
$+6.42\%$ / $8$-of-$8$ with a paired gain of $-0.186$ AA $[-0.220, -0.150]$. Resolving the
triple per origin instead (\S\ref{sec:gatesel}) moves the mean by $+0.0124$ AA, with a
bootstrap interval of $[-0.0735, +0.0865]$ that is an order of magnitude wider than the
point estimate. The two arms are not distinguishable, and the correction was made because
the earlier protocol was indefensible, not because the number demanded it. Two things the
correction does \emph{not} settle: it widens the paired interval on the kernel's gain from
$[-0.220,-0.150]$ to $[-0.272,-0.084]$, and it costs one transition against the parent.

\item \textbf{$\mu$ is no longer pinned, but the grid ceiling still binds sometimes.}
Under the superseded arm the selector chose $\mu = 0.4$ --- the largest value offered ---
at every origin, so the optimum sat on a boundary and the reported gain was a lower bound.
Under the causal arm it chooses $0.4$ at four of the eight scored origins and the interior
values $0.05$, $0.1$, $0.2$ at the rest. The boundary therefore still binds at half the
panel and the grid $\{0, 0.05, 0.1, 0.2, 0.4\}$ should still be extended, but the defect is
now partial rather than total. The change is unsurprising: a differently tempered kernel
wants a different mixing weight.

\item \textbf{The backbone is a development set, and the kernel has now partly
replicated.} The model was designed and tuned across many rounds on exactly these
transitions. In the round that produced the first version of this result, three other
statistically significant backbone findings were checked on the out-of-sample Huddleston
2020 panel and none replicated; one reversed sign and became significantly harmful. The
kernel has since been checked there and does replicate on the 23 validation origins
($-0.1455$ AA $[-0.206, -0.079]$, $21$ of $23$; \S\ref{sec:result}). The eight
\emph{test} origins of that panel --- the strictest split available --- are still running,
and until they report the out-of-panel evidence is one arm rather than two.

\item \textbf{The out-of-panel replication carries a transferred constant.} That H2020 run
used $(\tau_e, \tau_k, m) = (0.25, 0.02, 0)$, the operating point the \emph{backbone} gate
chose. The backbone gate read targets from 2018--2023 and every H2020 origin predates them,
so for that panel these are future information --- exactly the defect \S\ref{sec:gatesel}
removes on the backbone, left in place there because the run was most of the way through
when it was found. An H2020 gate table would fix it and has not been built.

\item \textbf{The comparison protocol is now symmetric, and that is worth stating in both
directions.} The rows in \S\ref{sec:result} attributed to Huddleston 2020 are scored with
their hyperparameters re-chosen at every origin on our folds, not held at the config their
papers ship, because comparing a per-origin selector against constants frozen in 2019
compares protocols as much as models. Averaged over their six models this \emph{costs}
them $0.042$ AA --- three folds over $16$ to $144$ candidates mostly adds variance to a
one- or two-coefficient model --- but for \texttt{lbi} it is worth $+0.467$ AA, taking it
from $8.144$ / $-2.56\%$ / $2$-of-$8$ to $7.677$ / $+4.20\%$ / $5$-of-$8$. The choice moves
\texttt{mlr+mom} from fourth to third of six fitted models, so it is not a choice that
flatters this work.

\item \textbf{The estimator and the score are different functionals.} Coefficients
minimise penalised cross-entropy against the assignment target \eqref{eq:assign}; every
number reported is an earth-mover's distance \eqref{eq:emd} to the realised season. The
gap between the two is not quantified here.

\item \textbf{$\delta$ contributes little.} The matched control \texttt{kcons}
($\delta = 0$) scores $-0.174$ against \texttt{kcmom}'s $-0.182$. The level term does
essentially all the work and the momentum extrapolation adds $0.008$ AA, which falsifies
the motivating argument for the arm even though the arm itself succeeds. Both figures are
from the superseded pooled arm; \texttt{kcons} has not been rerun under
\S\ref{sec:gatesel}, so the comparison is between two models at a common operating point
rather than each at its own.
\end{enumerate}

\end{document}
